% language: swedish
% figure: fig1 0.33 0.495 0.54 0.53
% figure: fig2 0.035 0.714 0.17 0.75
% figure: fig3 0.51 0.905 0.68 0.95
\textbf{Proposition 9.11 a:} Om $z_0$ hävbar singularitet då $\operatorname{Res}_{z_0} f = 0$

\textbf{Proposition 1.3 R:} Om $f(z) = \dfrac{g(z)}{(z - z_0)^m}$, $g$ holomorf, då $\operatorname{Res}_{z_0} f = \dfrac{g^{(m-1)}(z_0)}{(m - 1)!}$

\textbf{Proposition 9.14:} Om $f, g$ holomorf, $z_0$ enkelt nollställe till $g$, då $\operatorname{Res}_{z_0} \dfrac{f}{g} = \dfrac{f(z_0)}{g'(z_0)}$

\subsection{Tenta 2013 oktober uppg.\ 2}
Beräkna $\displaystyle\int_0^\pi \frac{dt}{1 + \sin^2 t}$

\textbf{Lösning:} Vi observerar $\sin^2(t) = \sin^2(t + \pi) \Rightarrow \displaystyle\int_0^\pi \frac{1}{1 + \sin^2 t}\,dt = \frac{1}{2} \int_0^{2\pi} \frac{1}{1 + \sin^2 t}$

Variabelbyte $z = e^{it}$, $dt = \dots = -\dfrac{i\,dz}{z}$
\[ \sin t = \dots = \frac{z^2 - 1}{2iz},\ 1 + \sin^2 t = \dots = -\frac{z^4 - 6z^2 + 1}{4z^2},\ \frac{dt}{1 + \sin^2 t} = \dots = f(z)\,dz \]
\[ f(z) = \frac{4iz}{z^4 - 6z^2 + 1}, \quad z^4 - 6z^2 + 1 = \prod_{j = 1}^4 (z - z_j),\ z_1 = \sqrt{3 + 2\sqrt{2}},\ z_2 = -\sqrt{3 + 2\sqrt{2}} \]
\[ \textcircled{\scriptsize $z_3$} = \sqrt{3 - 2\sqrt{2}} \quad \textcircled{\scriptsize $z_4$} = -\sqrt{3 - 2\sqrt{2}} \]
\textcolor{red!70!black}{\textbf{OBS!}} $0 < 3 - 2\sqrt{2} < 1$
\notefigure{fig1}{}
\[ \operatorname{Res}_{z_3} f = \left( \frac{4iz}{4z^3 - 12z} \right)_{z = z_3} = \dots = \rattat{-\frac{i}{2\sqrt{2}}}{-\frac{1}{2\sqrt{2}}}, \quad \operatorname{Res}_{z_4} f = \left( \frac{4iz}{4z^3 - 12z} \right)_{z = z_4} = -\frac{i}{2\sqrt{2}} \]
\[ \Longrightarrow \int_0^\pi \frac{1}{1 + \sin^2 t}\,dt = \frac{1}{2} \int_0^{2\pi} \frac{1}{1 + \sin^2 t}\,dt = \frac{1}{2} \int_{|z| = 1} f\,dz \underset{\textcolor{magenta!80!black}{\text{residysatsen}}}{=} \frac{1}{2} 2\pi i\left( -\frac{i}{2\sqrt{2}} - \rattat{\frac{i}{2\sqrt{2}}}{\frac{i}{\sqrt{2}}} \right) = \frac{\pi}{\sqrt{2}} \]

\noindent\rule{\linewidth}{0.4pt}

\begin{definition}
$f$ \emph{meromorf} om $f$ holomorf förutom poler.
\end{definition}
\begin{definition}
Om $\gamma$ sluten kurva i $\mathbb{C} \setminus \{0\}$ så definierar vi \emph{vindningstalet} $W(\gamma)$ som är \# varv $\gamma$ går runt $0$ i positiv riktning
\notefigure{fig2}{}
\end{definition}

\begin{theorem}[Argumentprincipen (Sats 3.2R)]
Antag $f$ meromorf i $G$, $\gamma$ styckvis glatt sluten enkel positivt orienterad nollhomotop kurva i $G$ som undviker $f$:s nollställen \& poler.

Då är $W(f \circ \gamma) = N(f, \gamma) - P(f, \gamma)$, där $N(f, \gamma) = $ \# nollställen till $f$ i $\operatorname{int}(\gamma)$, $P(f, \gamma) = $ \# poler till $f$ i $\operatorname{int}(\gamma)$
\notefigure{fig3}{}
\end{theorem}
